\documentclass[10pt,a4paper]{article}
\usepackage[a4paper,margin=1.35cm,top=1.1cm,bottom=1.1cm]{geometry}
\usepackage[T1]{fontenc}
\usepackage{lmodern}
\usepackage{microtype}
\usepackage{amsmath}
\usepackage{booktabs}
\usepackage{tabularx}
\usepackage{enumitem}
\usepackage{titlesec}
\setlist{nosep,leftmargin=1.1em}
\titleformat{\section}{\normalsize\bfseries}{\thesection}{0.5em}{}
\titleformat{\subsection}{\small\bfseries}{\thesubsection}{0.5em}{}
\titlespacing*{\section}{0pt}{5pt}{1.5pt}
\titlespacing*{\subsection}{0pt}{3pt}{1pt}
\setlength{\parindent}{0pt}
\setlength{\parskip}{1.5pt}
\newcolumntype{Y}{>{\raggedright\arraybackslash}X}
\pagestyle{empty}

\begin{document}

{\large\bfseries Ball possession from automatic tracking: stoppages first, then persistent control}\\[1pt]
{\small More than a Hack 2027, Phase~1 (FC~Barcelona vs Athletic Club). The implemented method, not planned work.}

\section{Inputs, training and validation data}
The tactical automatic tracking (players and ball) defines the timeline: submission frame $i$ is tracking record $i{+}1$. Also
used: the published phase table (halves, graded frames, kickoffs), the broadcast tracking and homography (aligned in Sec.~3), and
the \emph{set-piece} intervals of the automatic Smart Tagging, as fallible corroboration only. Video pixels are not used.
\textbf{No model is trained on labels.} The optional team-shape stage (Sec.~7, not used for the submission) is fitted per match
on the pipeline's own unambiguous frames, i.e.\ algorithmic pseudo-labels. \textbf{Validation data:} Metrica's three open sample
matches (anonymised; same provider, formats and 25\,fps; events annotated by Metrica, about 95 graded minutes each), used only to
compare settings (Sec.~9). No human labels, events or corrected tracking of the challenge match are used; its public labels
served only the final check.

\section{Pipeline overview}
One command runs these stages in order; all are deterministic and none reads a label.
\begin{center}\small\vspace{-4pt}
\begin{tabularx}{\linewidth}{@{}l>{\raggedright\arraybackslash}p{0.47\linewidth}Y@{}}
\toprule
stage & what it does & output \\
\midrule
1 load and check & parse both tracking exports; enforce frame order, frame rate and the phase table; content-keyed caches & players and ball per frame, in metres \\
2 synchronize & one broadcast offset per half, from player layouts (Sec.~3) & offsets, quality figures \\
3 clean players & bridge lost tracks, complete missing players from the live broadcast, one goalkeeper slot per team & players per frame \\
4 ball & drop impossible jumps, bridge gaps up to 5\,s, fill longer gaps from the live broadcast; also keep the tactical ball alone & two ball tracks \\
5 stoppages & judge every long tactical ball gap, add other dead-ball evidence (Sec.~4) & dead / live per frame \\
6 restarts & team of each restart, keeper holds (Sec.~5) & restart team per live stretch \\
7 possession & control evidence and the possession state machine (Sec.~6) & H, A, D or X per frame \\
8 (optional) team shape & veto and fill (Sec.~7); off for the submission & revised labels \\
9 write and check & official validation before and after writing; report with input/code hashes & submission, report \\
\bottomrule
\end{tabularx}\vspace{-4pt}
\end{center}

\section{Preprocessing and synchronization}
Every tracking frame must occur once and in order; incomplete coordinates count as missing; positions are in metres
($105.3\times68$\,m). Lost player tracks are bridged for up to 10\,s if the player could have run the distance (7\,m/s plus 2\,m
slack); missing players are completed from the broadcast where its player layout matches the tactical one; each team keeps one
goalkeeper slot (the marked keeper track, else the clearly deepest player in his own box, else his last position, which is never
used as evidence). Single ball detections implying jumps above 40\,m/s are discarded; ball gaps up to 5\,s are bridged linearly,
for continuity only (an interpolated ball never establishes control). \emph{Synchronization:} one frame offset per half, from
400 evenly spaced broadcast frames with a valid homography (reprojection error at most 3) and at least 8 players, each matched to
the tactical frame whose same-team player positions fit best over 0, 1 and 2\,s; accepted only with at least 10 matches, 85\% of
them within 0.5\,s of the median offset, and a clock slope within $5\cdot10^{-4}$ of 1 (obtained: 99.5\%/98.4\%, slope 0.99998).
The broadcast ball fills tactical gaps longer than 5\,s only where the broadcast shows the moment live: valid homography and at
least 4 players matching the tactical ones within 3\,m RMS (replays fail this check).

\section{Dead ball (D)}
\emph{Label-free observation.} The tactical ball track has gaps either below 1\,s or of 5\,s and more. The long gaps cover about
40\% of the graded frames, the usual out-of-play share; during them the players' median speed is about 1\,m/s (87\% of the
frames at walking pace, against 1.6\,m/s with the ball visible), the broadcast usually cuts away, and tagged set pieces begin
shortly before they end. Metrica's labelled sample matches show the same convention: the ball is present in 91--97\% of live
frames but in 0--8\% of dead ones. \emph{Rules.} Each tactical ball gap of 5\,s or more is judged on the tactical ball itself,
before any broadcast filling (a ball glimpsed by the broadcast during a stoppage must not hide it):
\begin{itemize}
\item \textbf{goalkeeper hold} (live): the same goalkeeper controls the ball (as in Sec.~6) where it vanishes and where it
      reappears, and it was not last seen at a line; the ball is in play (Law~12) and his team possesses it;
\item \textbf{dead}: the ball was last seen out of play; or the players' median speed over the gap is at most walking pace
      (1.5\,m/s); or at most 2.0\,m/s with one further cue (last seen at a line, reappearing at a restart spot, at least 12
      players within 10\,m for 3\,s, or a tagged set piece whose interval contains the restart, $\pm$1\,s). A tag never decides alone;
\item otherwise \textbf{live}: play continues unseen, and the owner is kept.
\end{itemize}
Also dead: the observed ball more than 1.5\,m beyond a line for at least 0.4\,s; the wait before a detected placed-ball restart
(Sec.~5); graded frames before a half's kickoff (Law~8; the kickoff is where the published phase table starts the clock). On the
challenge match 109 of 130 long gaps are judged dead (92 by walking pace, 13 by the ball out, 4 corroborated), 6 are goalkeeper
holds and 15 stay live; 42\% of the graded frames are dead, in 143 stoppages.

\subsection{Bridging (implemented and tested, off for the submission)}
After the rules above, a live stretch can lie between two stoppages. \textbf{Bridging} would relabel such a stretch as dead when
it is shorter than 5\,s and the players' median speed over it stays at walking pace: the idea was that the tracker had only
glimpsed the ball during one long stoppage (a player picking it up for a throw-in) and play never restarted. \textbf{Without
bridging} the stretch stays live: it is a restart, gets a restart team (Sec.~5) and becomes a short possession. On the challenge
match bridging would join 29 such stretches (1{,}972 frames, 2.7\,s on average) into 114 instead of 143 stoppages. On the
labelled matches these stretches were mostly real play---a quick throw-in or free kick followed by the next stoppage, taken
while the players had not yet accelerated---and bridging lowered spellF1 by 0.008--0.012 in all four conditions, winning in
only 1 of 12 match--condition pairs (Sec.~9). It therefore stays available as a setting but is off.

\section{Restarts}
Each live stretch begins with a restart, whose team possesses the ball until the other team's control is confirmed:
\begin{enumerate}
\item a \emph{detected placed-ball restart}: the ball lies still (rolling median speed at most 1.5\,m/s) for at least 1\,s and is
      then played; its location gives the kind (corner within 2.5\,m of a flag, goal kick in a goal area, throw-in on a touchline,
      kickoff or penalty at the spot, else free kick); it is accepted if it lasts long enough for its kind (1--3\,s), the nearest
      opponent stands back (1.5\,m at a throw-in, 3--6\,m at the others; the laws ask 2 and 9.15\,m), the players are slow (at
      most 2\,m/s), the taker stands within 3\,m and the ball then moves. Its taker's team restarts; the still stretch is dead;
\item else the \emph{first ball seen within 1\,s} of the stretch start: at a corner flag the attacking team, in a goal area the
      defending team (attack directions from the teams' positions at each kickoff), elsewhere the team of the player in control
      at the ball (the taker);
\item else none: the first team whose control is confirmed in the stretch has it from the start.
\end{enumerate}
On the challenge match 143 live stretches start as: 18 placed-ball restarts, 6 corners, 12 goal kicks, 64 takers, 43 unknown.

\section{Possession (H/A)}
A frame carries \emph{control evidence} for a team when the observed (not interpolated) ball moves at most 12\,m/s (where its speed
is known) and one of the team's players is within 3\,m of it and more than 1\,m nearer than every opponent. Each live stretch is
decoded left to right:
\begin{enumerate}
\item the owner starts as the restart team (Sec.~5), or unknown;
\item a run of evidence for one team is counted frame by frame; any frame without it, or with the other team's, resets the count;
\item when the count reaches 0.6\,s (15 frames) for a team that is not the owner, possession changes; the change is dated back to
      the first frame of the run (the prediction is offline). If the owner was unknown, that team owns the stretch from its start;
\item otherwise the owner keeps the ball through passes, loose and contested balls, deflections, brief touches and occlusions of
      any length; a goalkeeper hold gives his team the ball from his catch;
\item a stoppage or the end of a half ends the stretch. No smoothing follows; X remains only in a stretch whose team is never
      established (230 of 158{,}857 graded frames; the official scorer repeats the previous label there).
\end{enumerate}

\section{Optional stage: team shape (off for the submission)}
The team in possession spreads out, advances and runs less than the team chasing it. Per frame, three home-minus-away
differences of the outfield players are formed: mean distance to the team centre, speed of the centre towards the opponent goal,
and median running speed. A logistic model of the team in possession given these differences is fitted per match on the frames
where the owning team's player has the ball at his feet (in-sample AUC about 0.88 on the challenge match). Shape follows a change
of possession only about 3\,s later, so it never times a change; it judges one. \textbf{Veto:} a possession won in open play that
lasts at least 5\,s, but whose part from 3\,s after its start the model gives its winner a probability below 0.2, is returned to
the previous owner. \textbf{Fill:} unattributed frames take the team the shape gives at least 0.8. On the challenge match it would
veto 7 possessions and change 0.8\% of the graded frames; on the validation matches it gave no consistent gain (Sec.~9).

\section{Recommended settings (submitted)}
These are the values in the configuration that produced the submitted file; ``measured'' means fixed from label-free statistics
of the challenge input, ``validated'' means compared on the labelled sample matches (Sec.~9).
\begin{center}\small\vspace{-4pt}
\begin{tabularx}{\linewidth}{@{}lrY@{}}
\toprule
setting & value & basis \\
\midrule
long ball gap (stoppage candidate; also the interpolation limit) & 5\,s & measured: no gaps between 1 and 5\,s \\
walking pace / slow pace (median player speed) & 1.5 / 2.0\,m/s & walking speed 1.4\,m/s; validated (1.2 and 1.8 change nothing) \\
ball beyond a line before it counts as out & 1.5\,m for 0.4\,s & tracking error about 1\,m; rejects one-frame glitches \\
bridging of short walking-pace stretches & off & validated (Sec.~4.1) \\
restart team window & 1\,s & validated: without restart teams spellF1 drops 0.04--0.05 \\
control radius / margin over the nearest opponent & 3\,m / 1\,m & reach plus tracking error; validated under realistic noise \\
maximum ball speed for control & 12\,m/s & a fast pass near a player is not control \\
confirmation of a change of possession & 0.6\,s & validated (0.4\,s before; 0.2 and 1.2\,s worse) \\
broadcast synchronization and completion & on & offsets computed and gated (Sec.~3) \\
Smart Tagging set pieces & on, $\pm$1\,s & corroboration only; decides no gap alone \\
team-shape stage & off & validated: no consistent gain \\
\bottomrule
\end{tabularx}\vspace{-4pt}
\end{center}

\section{Validation on other matches}
The unchanged pipeline runs on Metrica's three sample matches. Reference labels follow Metrica's event definitions: the team of
the latest on-ball event has the ball (\emph{touch}); dead from a ball out, goal, offside, foul or referee stop until the next
on-ball event. A \emph{controlled} variant applies the challenge's proxies: a spell won in open play counts only if it lasts 2\,s
or contains a completed pass. The sample tracking is corrected (the ball lies about 0.1\,m from its owner), so the benchmark is
also run with a smooth random ball error (1.15\,m per axis, 3\,s time constant), which reproduces the challenge's distances
between ball and nearest player (median 1.8\,m) and its ball speeds, both measured without labels. Scores use the official
metrics: spellF1, mean over the three matches, and in how many of the 12 match--condition pairs a variant beats the submitted
settings. \emph{Oracle} rows replace one component by the truth.
\begin{center}\small\vspace{-4pt}
\begin{tabular}{@{}lccccc@{}}
\toprule
variant (one change from the submitted settings) & clean, touch & clean, contr. & noisy, touch & noisy, contr. & better \\
\midrule
\textbf{submitted (0.6\,s, 3\,m, 1\,m, no bridging)} & \textbf{0.569} & \textbf{0.659} & \textbf{0.527} & \textbf{0.630} & -- \\
previous settings (0.4\,s, bridging) & 0.548 & 0.593 & 0.512 & 0.575 & 0 \\
confirmation 0.2 / 0.4\,s & 0.503 / 0.556 & 0.493 / 0.601 & 0.448 / 0.520 & 0.461 / 0.584 & 0 / 1 \\
confirmation 0.8 / 1.2\,s & 0.559 / 0.518 & 0.676 / 0.649 & 0.497 / 0.446 & 0.625 / 0.573 & 5 / 0 \\
control radius 1.5 / 2\,m & 0.545 / 0.561 & 0.673 / 0.673 & 0.414 / 0.465 & 0.531 / 0.583 & 3 / 3 \\
control radius 2.5 / 4\,m & 0.569 / 0.551 & 0.673 / 0.623 & 0.508 / 0.526 & 0.617 / 0.618 & 5 / 3 \\
margin 0.5 / 2\,m & 0.585 / 0.536 & 0.657 / 0.661 & 0.520 / 0.476 & 0.596 / 0.611 & 5 / 2 \\
walking pace 1.2 / 1.8\,m/s & 0.568 / 0.569 & 0.658 / 0.659 & 0.524 / 0.527 & 0.627 / 0.629 & 0 / 0 \\
bridging 5\,s & 0.560 & 0.651 & 0.517 & 0.618 & 1 \\
no restart teams & 0.525 & 0.620 & 0.485 & 0.585 & 0 \\
team shape (veto 0.2; 0.1 / 0.3 similar or worse) & 0.563 & 0.652 & 0.529 & 0.633 & 7 \\
\midrule
oracle dead ball / oracle team & 0.608 / 0.919 & 0.711 / 0.900 & 0.554 / 0.933 & 0.663 / 0.924 & 12 / 12 \\
\bottomrule
\end{tabular}\vspace{-4pt}
\end{center}
The submitted settings beat the previous ones in all 12 pairs. No single change beats them in more than 7 of 12 pairs: smaller
radii or margins and 0.8\,s win only in some clean-data conditions (1.5\,m would lose about 0.1 with realistic noise), and team
shape gains 0.002--0.003 with noise but loses 0.006--0.007 on clean data, too little and too mixed to switch it on. The oracles locate the remaining error: perfect stoppages would add 0.03--0.05, perfect team calls in
live play 0.24--0.41.

\section{Tests and reproduction}
187 synthetic tests (no licensed data needed) cover, by area:
\begin{center}\small\vspace{-4pt}
\begin{tabularx}{\linewidth}{@{}lrY@{}}
\toprule
area & tests & what is checked \\
\midrule
loading, frames, caches & 40 & frame IDs repeated, unordered or out of range; partial coordinates; phase table coverage; cache invalidated by content, parser or corruption; archive unpacking \\
synchronization & 12 & matching by player layout, offsets per half, replay rejection, quality gates (agreement, slope, matches) \\
players, keepers, ball & 34 & track bridging, broadcast completion, keeper placement, false detections, gap filling, speeds \\
stoppages and restarts & 45 & gap verdicts, goalkeeper hold, Smart Tagging never deciding alone, out-of-pitch runs, pre-kickoff, bridging on/off and never across halves, restart kinds and teams, judged on the tactical ball \\
possession model & 20 & retention through occlusion, confirmation dated back, restart and keeper teams, half resets, end to end \\
validation data, team shape & 13 & event-to-label rules (goal kick is no goal, ball lost, fouls), controlled labels, noise reproducibility and size, shape features, fit, veto and fill \\
viewers and helpers & 23 & plots, tables, progress output \\
\bottomrule
\end{tabularx}\vspace{-4pt}
\end{center}
\texttt{python run\_pipeline.py -{}-evaluate} (repository root) writes \texttt{outputs/current/submission.csv} (199{,}319 rows,
officially validated before and after writing) and a report hashing all inputs and sources, saved before any label is read.
The benchmark is \texttt{python -m src.benchmark -{}-variants} in \texttt{development/}. Warm, repeated and cold-cache runs
are byte-identical (SHA-256 \texttt{9664234e\ldots f956}); the official validator accepts the file.

\section{Public results, scope and limitations}
Measured once on the public five minutes (7{,}500 frames, 17 possessions) after the settings were fixed on the validation
matches; no change followed:
\begin{center}\small\vspace{-5pt}
\begin{tabular}{@{}lll@{}}
\toprule
 & value & bootstrap 5--95\% \\
\midrule
spellF1 (ranking metric) & 0.526 & 0.366--0.663 \\
macroF1 \quad(per-class F1: H 0.857, A 0.802, D 0.552) & 0.737 & 0.557--0.858 \\
possessions / mean duration \quad(truth: 17 / 14.2\,s) & 18 / 14.6\,s & \\
\bottomrule
\end{tabular}\vspace{-5pt}
\end{center}
The previous settings scored 0.544 and 0.739 here. Seventeen possessions cannot rank methods: the scorer treats spellF1 gaps below
about 0.30 as noise, and by the rules this window may not choose between versions; the choice rests on Sec.~9.
Most frame errors in the window are dead-ball frames labelled live (D recall 0.45).
\textbf{Limitations.} (1)~Dead ball comes from when the tracker loses the ball and how the players move: a stoppage starts late
while the ball stays in view, and unseen slow live play may be taken for one. (2)~Possession changes in live play are the main
error (Sec.~9): proximity cannot always tell a touch from control, and the best confirmation time depends on how control is
defined. (3)~The validation labels are event-based proxies on three matches with corrected tracking and simulated ball error.
(4)~Restart teams fail when the ball reappears in flight; one broadcast offset per half; Smart Tagging teams and phases are unused.

\end{document}
